\chapter{实验与结果分析}

本章对系统进行全面的实验评测，依次验证视觉无感性、攻击防御效果、目标检测与高级攻击防御、对抗噪声消融、运行性能，并诚实分析方案的局限性。所有结论均以固定的基准评测（benchmark）与语料实验结果为准，避免单次演示的偶然波动。

\section{实验环境与数据集}

实验在 macOS 平台、Python 3.10 环境下进行，主要依赖 NumPy、SciPy、PyCryptodome、OpenCV、moderngl，OCR 评测使用 Tesseract 5.4.1、EasyOCR 1.7.2、PaddleOCR 3.6.0，目标检测使用 Ultralytics YOLOv8n（8.4.65）。系统配套 176 项单元测试，运行 \verb|pytest tests/ -q| 全部通过，保证被测算法实现正确。

评测数据集为自建的 120 样本文本语料，由 12 类模板各生成 10 个确定性变体，覆盖英文句子、中文句子、数字、账号口令、URL 令牌、代码片段、表格、段落、财务字段等敏感信息类型。每个样本附带真实文本（ground truth）与类别、语言、版式、字号等元数据，用于计算字符识别准确率并支持分层统计。OCR 结果均在 Tesseract、EasyOCR、PaddleOCR 三个引擎上跑满全量语料，用于验证防御效果不依赖单一 OCR 引擎。系统另实现了在线 VLM 可读性评测入口，可将攻击帧发送至 OpenAI-compatible 视觉语言模型接口并在本地计算同一套恢复指标；由于本章未完成真实在线调用，后续不引用 VLM 数值。

\section{视觉无感性评估}

视觉无感性通过闪烁感知指数（FPI）、CIEDE2000 色差（$\Delta E$）等客观指标量化。图5-1展示了 $n=4$ 配置下的视觉积分对比：第一行为原始图像、人眼积分还原图与放大 4 倍的差异图，第二行为相机视角下的各个单子帧。可见人眼积分还原图与原图几乎无差异，而单子帧均为无法辨识的稀疏点阵。

\begin{figure}[htbp]
  \centering
  \includegraphics[width=0.92\textwidth]{demo_n4_grid.png}
  \caption{视觉积分对比（上：原图/人眼积分/差异；下：相机视角单子帧）}
\end{figure}

不同参数配置下的视觉无感指标如表5-1所示，数据取自固定基准评测。可见在满足 $f_r \ge 60n$ 的配置（$n=2/144$Hz、$n=4/240$Hz）下，FPI 均远低于 0.1 的人眼可感知阈值；而不满足约束的配置（$n=4/144$Hz）FPI 高达 0.30，验证了刷新率约束的必要性。积分还原帧相对原图的色差 $\Delta E$ 在 $n=4$ 时约为 0.37（该值为含对抗噪声路径、跨语料样本的平均，受测试图像与噪声预算影响），远低于 1.0 的不可感知阈值。

\begin{table}[htbp]
  \centering
  \caption{不同配置下的视觉无感与单帧防御指标}
  \begin{tabular}{ccccc}
    \toprule
    配置 & FPI & $\Delta E$ & 单帧互信息比 & 单帧 OCR \\
    \midrule
    $n=2$, 60Hz   & 0.250 & 0.19 & 0.612 & 0.0\% \\
    $n=2$, 144Hz  & 0.014 & 0.19 & 0.612 & 0.0\% \\
    $n=4$, 144Hz  & 0.300 & 0.37 & 0.342 & 0.0\% \\
    $n=4$, 240Hz  & \textbf{0.030} & \textbf{0.37} & \textbf{0.343} & \textbf{0.0\%} \\
    \bottomrule
  \end{tabular}
\end{table}

需要说明的是，"单帧互信息比"（单子帧相对原图的归一化互信息）在 $n=4$ 时约为 0.34，该指标对测试图像内容、噪声预算与统计口径较为敏感（不同样本与运行可能在 0.2 至 0.6 之间波动），因此本文将其作为辅助参考，而以字符识别准确率（OCR）作为衡量隐私保护的主指标。

\section{攻击防御实验}

\subsection{语料级 OCR 防御效果}

使用扩展语料对原始图像与单子帧的字符识别准确率进行评测，结果如表5-2所示。Tesseract、EasyOCR、PaddleOCR 在 120 样本全量语料上的原始帧识别准确率分别为 94.0\%$\pm$9.1\%、94.1\%$\pm$8.1\%、94.9\%$\pm$7.3\%，而单子帧识别准确率分别降至 0.0\%$\pm$0.2\%、0.0\%$\pm$0.0\%、0.0\%$\pm$0.0\%；三者配对降幅均达到 93.9\% 以上。除字符准确率外，本文进一步统计词级准确率、整句精确匹配率和敏感 token 恢复率。单子帧在三种 OCR 引擎下的词级准确率、exact-match 与敏感 token 恢复率均为 0.0\%，其中敏感 token 统计覆盖 104 个包含账号、数字、URL、代码或凭据字段的样本。分层统计显示，代码、财务字段、数字、段落、表格、URL 令牌等类别的单子帧识别率均接近 0，其中表格类因原始 OCR 本身较难，原始帧准确率相对较低。上述结果表明，时间分片对文本信息的防护并非只对某一 OCR 引擎有效。

\begin{table}[htbp]
  \centering
  \caption{语料级 OCR 防御效果（$n=4$）}
  \begin{tabular}{ccccc}
    \toprule
    OCR 引擎 & 样本数 & 原始图像准确率 & 单子帧准确率 & 防御效果 \\
    \midrule
    Tesseract & 120 & 94.0\%\,$\pm$\,9.1\% & 0.0\%\,$\pm$\,0.2\% & 有效 \\
    EasyOCR   & 120 & 94.1\%\,$\pm$\,8.1\% & 0.0\%\,$\pm$\,0.0\% & 有效 \\
    PaddleOCR & 120 & 94.9\%\,$\pm$\,7.3\% & 0.0\%\,$\pm$\,0.0\% & 有效 \\
    \bottomrule
  \end{tabular}
\end{table}

\subsection{相机采集攻击防御}

针对不同相机采集策略，使用相机模拟器进行攻防实验，结果如表5-3所示。对于目标威胁模型内的全局快门抽帧、流式逐帧识别、卷帘快门混合采样，以及不足一个完整周期的时域平均，单帧 OCR 准确率均为 0\%；对于长曝光攻击，在未插入反色帧时攻击者可积分还原（100\%），而插入反色帧后被有效对抗（0\%），验证了反色帧防御的有效性。

\begin{table}[htbp]
  \centering
  \caption{相机采集攻击防御效果}
  \begin{tabular}{ccc}
    \toprule
    攻击方式 & OCR 准确率 & 是否在目标威胁模型内 \\
    \midrule
    原始帧（基线） & 100\% & — \\
    全局快门单帧抽帧 & \textbf{0\%} & 是（有效防御） \\
    流式逐帧识别 & \textbf{0\%} & 是（有效防御） \\
    卷帘快门混合采样 & \textbf{0\%} & 是（有效防御） \\
    时域平均（$<1$ 周期） & \textbf{0\%} & 是（有效防御） \\
    长曝光（插入反色帧） & \textbf{0\%} & 是（有效防御） \\
    长曝光（无反色帧） & 100\% & 否 \\
    时域平均（$\ge 1$ 完整周期） & 100\% & 否（已知局限，见 5.6） \\
    \bottomrule
  \end{tabular}
\end{table}

\section{屏幕-相机通信式强攻击者评估}

前述实验主要覆盖被动抽帧、逐帧流式识别、卷帘快门混合采样和长曝光等常见采集策略。为进一步贴近投稿论文所需的强攻击者模型，本文借鉴屏幕-相机通信与亮度隐蔽信道研究中的解码思路，补充构造了三类更主动的连续视频攻击：相位搜索重构、加权差分累加和单通道重构。该类攻击者不假设已知显示周期起点，但知道屏幕存在时域调制，会在连续帧中搜索候选周期窗口，并检查亮度通道或蓝色通道是否存在可恢复的残差信息。

单样本攻击结果如表5-4所示。相位搜索均值重构和相位搜索最大值重构均可将 OCR 准确率恢复至 100\%，说明攻击者即使不知道周期起点，也能通过连续视频搜索找到可重构窗口；蓝通道最大值重构同样恢复至 100\%，提示颜色通道不应被假设为安全侧信道。相比之下，加权差分累加在该样本上未恢复 OCR，说明单纯累加相邻帧变化不足以直接恢复文本。

\begin{table}[htbp]
  \centering
  \caption{屏幕-相机通信式强攻击者单样本评估（$n=4$，Tesseract）}
  \begin{tabular}{ccc}
    \toprule
    攻击方式 & OCR 准确率 & 结论 \\
    \midrule
    相位搜索均值重构 & 100.0\% & 失守 \\
    相位搜索最大值重构 & 100.0\% & 失守 \\
    加权差分累加（亮度通道） & 0.0\% & 当前样本有效防御 \\
    加权差分累加（蓝通道） & 0.0\% & 当前样本有效防御 \\
    蓝通道最大值重构 & 100.0\% & 失守 \\
    \bottomrule
  \end{tabular}
\end{table}

为避免仅由单样本得出结论，本文进一步在 120 样本文本语料上运行同一强攻击套件。实验采用 $n=4$、$\epsilon=8/255$、连续 2 个显示周期、Tesseract OCR，单次 OCR 超时设为 4s；结果如表5-5所示。全局快门单帧与 luma/blue 加权差分累加的字符恢复率仍为 0\%，但完整周期平均、相位搜索和蓝通道最大值重构在语料级均恢复出高比例文本。逐样本选择恢复率最高的攻击时，字符恢复率均值为 95.2\%（95\%CI 93.8\%--96.5\%），exact-match 为 91.7\%，敏感 token 恢复率为 98.6\%，且 120 个样本的字符恢复率均超过 20\% 泄露阈值。

\begin{table}[htbp]
  \centering
  \caption{120 样本语料上的强相机攻击统计（Tesseract）}
  \small
  \begin{tabular}{lccc}
    \toprule
    攻击方式 & 字符恢复率 & exact-match & 敏感 token 恢复率 \\
    \midrule
    全局快门单帧 & 0.0\% & 0.0\% & 0.0\% \\
    加权差分累加（luma/blue） & 0.0\% & 0.0\% & 0.0\% \\
    完整周期平均 & 94.3\% & 85.8\% & 98.4\% \\
    相位搜索均值重构 & 94.2\% & 87.5\% & 97.2\% \\
    相位搜索最大值重构 & 94.3\% & 85.8\% & 96.4\% \\
    蓝通道最大值重构 & 95.0\% & 90.0\% & 98.3\% \\
    最强攻击逐样本择优 & \textbf{95.2\%} & \textbf{91.7\%} & \textbf{98.6\%} \\
    \bottomrule
  \end{tabular}
\end{table}

该补充实验进一步收紧了本文的安全声明：系统能够有效防御被动抽帧、流式逐帧识别和短曝光采样，但不应宣称防御主动连续录制后的完整周期重构、相位搜索或单通道最大值重构。这一结果与本文关于线性时间分片方案边界的分析一致，也为后续引入非线性时域调制、内容自适应扰动和真实屏幕-相机闭环实验提供了明确方向。

\section{目标检测与高级攻击防御}

\textbf{目标检测防御。}使用 YOLOv8n 评估时间分片对物体级信息的防护。实验表明，单子帧的检测 mAP50 为 0.40（召回 0.40），明显高于单子帧 OCR 的 0\%。这说明目标检测依赖粗粒度的形状与颜色团块，对 $1/n$ 稀疏采样比依赖细笔画的文本识别更鲁棒。因此本文区分"文本/细结构信息"（防护强）与"物体级场景信息"（防护较弱）两类威胁，不笼统宣称等效防御。

\textbf{离轴相机攻击。}针对技术方案 7.2 节提出的视角差异化掩模，构造 35° 离轴、$3\times3$ 区域差异化掩模的攻击实验。结果显示：正视、离轴朴素捕获、离轴加强攻击者校正三种路径下，整周期平均还原的 OCR 准确率均为 100\%，离轴 SSIM 仅下降 0.064。其原因在于整周期平均时每个区域无论采用何种差异化掩模都各自满足完备性、各自还原，掩模差异化并不参与防御。这一诚实结论表明：视角差异化掩模在纯软件可逆模型下不构成 OCR 级保护，真正的视角防护须依赖不可逆的物理 LCD 视角响应，属本原型范围外。

\textbf{学习型重构攻击。}训练一个轻量 U-Net 尝试从单子帧重构原图，作为学习型攻击的下界，其重构结果 OCR 准确率为 0\%。进一步的 inpainting 消融揭示了对抗噪声的关键作用：对不含噪声的"干净子帧"做单帧 inpainting 可将高对比文档 OCR 还原至 100\%，而对叠加了对抗噪声基底（pedestal）的子帧做同样攻击，OCR 跌回 0\%。这说明对抗噪声并非可有可无的边际增强，而是单帧防御对 inpainting 攻击保持鲁棒的关键支撑。该实验也修正了仅以 SSIM 度量的误判——SSIM 受白底背景主导，不能作为隐私指标，OCR 才是。

\section{对抗噪声消融实验}

为厘清掩模与噪声各自的贡献，进行四组消融对照。在默认 $n=4$、$\epsilon=8/255$ 场景下，随机点阵掩模通常已将单帧 OCR 压到 0，触发"地板效应"，掩盖了噪声的边际贡献；为此另设"弱掩模泄露"压力场景（模拟面板串扰、相机去噪使被关闭像素仍泄露部分亮度），此时 OCR 对比度目标噪声能提供可测的额外防御。综合消融结果可得出与上一节一致的结论：\textbf{掩模是主防御，对抗噪声是二级强化}，且在掩模泄露或面对单帧 inpainting 攻击时，噪声的作用不可或缺。

\section{运行性能实测}

本原型为 Python + NumPy 路径，性能实测结果如下表所示（含 GPU 不可用时的软件渲染回退）。在 1080p、$n=4$ 配置下，掩模生成约 15.5ms、端到端单周期生成约 451ms；1440p 下分别约 28.4ms 与 873ms。

\begin{table}[htbp]
  \centering
  \caption{运行性能实测（$n=4$）与技术方案声称值对比}
  \begin{tabular}{cccc}
    \toprule
    评估项 & 1080p 实测 & 1440p 实测 & 技术方案声称值 \\
    \midrule
    掩模生成 & 15.5ms & 28.4ms & $<0.1$ms（GPU CS） \\
    对抗噪声生成 & 355.6ms & 704.0ms & 查表法 \\
    子帧合成（SDR$\times4$） & 80.4ms & 140.5ms & — \\
    渲染单子帧 & 22--31ms & 40--55ms & 帧生成 $+0.3$ms \\
    端到端单周期 & 451.5ms & 873.0ms & 周期 16.7ms 预算 \\
    \bottomrule
  \end{tabular}
\end{table}

实测耗时显著高于技术方案声称的驱动层 GPU 指标（如 $<0.1$ms 掩模生成、$+0.3$ms 帧生成、GPU 负载 3--5\%）。这是因为技术方案的低开销指标假设核心算子以 GPU Compute Shader 并行执行并下沉至驱动/合成器层，而本原型运行于用户态 Python/CPU 路径。因此，\textbf{本文不以这些性能数字作为系统能力声明}；本原型的价值在于验证算法的正确性与防御的有效性，而非生产级实时性能。要达到技术方案的性能指标，需将核心算子移植至 GPU 驱动层，这属于产品化的未来工作。

\section{局限性分析}

本节诚实分析方案最重要的局限，这也是本课题最具价值的科学结论。如表5-3所示，当攻击者叠加覆盖一个完整显示周期（$\ge n$ 帧）并做时域平均时，可将图像完全还原，OCR 准确率回到 100\%。

其数学根因已在第三章给出：人眼"视觉无感"依赖完备性约束 $\sum_k I_k = I$ 的线性积分，而攻击者的时域平均使用同一个求和运算；互斥完备掩模在覆盖完整周期后被精确平均还原，对抗噪声因 $\sum_k N_k = 0$ 同时被抵消。\textbf{凡是能被人眼积分还原的，就能被覆盖完整周期的相机积分还原}——这是线性时域方案的原理性张力，而非实现缺陷。

因此，本系统的安全声明严格限定在目标威胁模型内：有效防御被动抽帧、流式识别、卷帘快门以及插入反色帧后的长曝光，这些覆盖了智能眼镜偷拍的真实采集策略；而对"主动连续录制完整周期 + 精确帧对齐 + 时域平均"这一更强的离线攻击，本文不予宣称防御。可行的缓解方向包括：利用物理 LCD 视角响应的不可逆视角差异化、引入人眼与相机积分差异的非线性时域调制、随机化周期长度与相位以增加完整周期对齐难度等，均列为未来工作。

\section{本章小结}

本章通过系统实验验证了：在 $n=4$、240Hz 配置下 FPI 为 0.030、$\Delta E$ 约 0.37，人眼视觉无感；Tesseract、EasyOCR、PaddleOCR 在 120 样本语料上的单子帧识别准确率均约为 0.0\%，配对降幅达到 93.9\%--94.9\%，且词级准确率、exact-match 与敏感 token 恢复率均为 0.0\%；对目标威胁模型内的各类相机攻击防御有效；对抗噪声是单帧防御对 inpainting 攻击鲁棒的关键；性能实测如实反映了 Python 原型与驱动级实现的差距。本章同时诚实揭示了完整周期对齐平均、相位搜索和单通道最大值重构这一类线性方案边界，其中 120 样本逐样本最强攻击字符恢复率为 95.2\%，并明确了系统的防护范围与缓解方向。
